\scanpage{24}
$\widetilde N$, c'est un $\underline O_X$-$G$-module \deleted{\unsure{sur $X$} \unsure{conte|contenu}}
$\subset\inserted{\unsure{\widetilde{\underline{\Gamma}^{\cdot}(F')}|\underline{\Gamma}^{\cdot}(F')}=}\deleted{\unsure{E}}F'$, et contenant $\widetilde{\underline{\Gamma}^{\cdot}(F)}=F$. C'est donc le
\deleted{\unsure{$F'$} cherché} faisceau cherché. Par \unsure{construct$^{\mathrm n}$|construction},
sa formation commute à \unsure{toute} extension
plate $S'\to S$ de la base. Enfin, si
$F$ est de type fini, \deleted{\unsure{alors $\exists$}} prouvons que
son enveloppe $F'$ l'est. On peut supposer $S$
affine. Alors $\exists\,n$ tel que $F(n)$
\deleted{\unsure{pour} \unsure{tt} $\unsure{i\leq N}$}
engendré par ses sections, \deleted{\unsure{soient des}} \inserted{donc} par des
sections $\varphi_i\deleted{\unsure{*}}\ (1\leq i\leq N)\in\Gamma(\unsure{\mathcal O_S|S},\underline{\Gamma}^n(F))$.

Considérons alors le sous-module
\deleted{$\underline N^n$ de $\underline{\Gamma}^{\unsure{\cdot|\bullet}}(E)$ engendré par les $\varphi_i$.
Alors $\widetilde{\underline N}^{\cdot}\subset F$ est en \deleted{\unsure{un}} fait égal à $F$,
comme on constate tt de suite. On
peut alors répéter la construction précédente,}
\unsure{ou|au} $G$-invariant de $\underline{\Gamma}^{\unsure{n|-}}(F)$ engendré
par les $\varphi_i$, soit $N^n$, et le sous
Module de $\underline E$ qu'il \unsure{définit|défint}. \uline{On
vérifie} \unsure{quand même} que c'est $\underline F'$, ce qui
prouve que $F'$ est de type fini\unsure{.....|......}
% slanted
\margin{Rédiger mieux \unsure{en|au|ce} début de la démonstration, pour pouvoir y référer ici \unsure{.....|...}}
% archive sheet number: 23
